\section{Applications and related topics}\label{sec-appl}



\subsection{Random involutions and random signed involutions}

The ensemble $S^\symmO_{n,m}$ is the set of involutions with 
$n$ 2-cycles and $m$ 1-cycles.
In the previous section, we considered the limiting statistics when 
$n$ and $m$ are related by $m=[\sqrt{2n}\alpha]$ with $\alpha$ 
being finite ; either fixed or $\alpha\to 1$ with certain rate.
It is of interest to consider the whole set of involutions 
without constraints on the number of fixed points.
Similarly, the signed involutions without constraint 
on the number of fixed points and negated points is also of interest.
We define the ensembles of involutions and signed involutions 
\begin{align}
  \tilde{S}_n &=\{\pi\in S_n : \pi=\pi^{-1} \}, \\
  \tilde{S}_n^\symmu &= \{\pi\in S_{2n} : 
\pi=\pi^{-1},\ \pi=\iota\pi\iota \}, 
\end{align}
and denote by $\tilde{L}_n(\pi)$ and $\tilde{L}^\symmu_n(\pi)$ the length
of the longest increasing subsequence of $\pi\in\tilde{S}_n$ 
and that of $\pi\in\tilde{S}_n^\symmu$, 
respectively.

Noting $\tilde{S}_n=\bigcup_{2k+m=n} S^\symmO_{k,m}$, 
we have 
\begin{equation}\label{e-appl-1}
  \Prob\bigl(\tilde{L}_n\le l\bigr)= \frac{1}{|\tilde{S}_{k,m}|} 
\sum_{2k+m=n} \Prob\bigl(L^\symmO_{k,m}\le l\bigr) |S_{k,m}|. 
%\quad p^\symmO_{kml}:=\Prob\bigl(L^\symmO_{k,m}\le l\bigr).
\end{equation}
It is not difficult to check that (see pp.66-67 of \cite{Kn}) 
as $n\to\infty$, the main contribution to the sum 
$|\tilde{S}_{k,m}|=\sum_{2k+m=n} |S_{k,m}|$ comes from 
$\sqrt{2k}-(2k)^{\epsilon+1/4}\le m\le\sqrt{2k}+(2k)^{\epsilon+1/4}$.
Comparing with the scaling $m=[\sqrt{2k}-2w(2k)^{1/3}]$ 
in Theorem \ref{thm3}, the main contribution to the sum 
\eqref{e-appl-1} comes from when $w=0$, or $\alpha=1$.
Thus we obtain GOE Tracy-Widom distribution function in the limit.
%An alternating way to see the above result is to note that 
%the length of the longest 
%increasing subsequence and that of the longest decreasing subsequence 
%of random involutions have the same statistics.
%This can be seen from the fact that the set of involutions 
%of $n$ letters is bijective to the set of standard Young tableaux 
%of size $n$ (see e.g. \cite{Kn}).
%By taking transpose of standard Young tableaux, 
%the first row and the first column have the same statistics 
%(the uniform probability distribution on the set of Young tableaux 
%has symmetry for rows and columns due to the hook formula).
Similarly, we can obtain the convergence of moments.
The signed involution case is analogous.

\begin{thm}
For fixed $x\in\R$, 
\begin{align}
  \lim_{n\to\infty} 
\Prob\biggl( \frac{\tilde{L}_n-2\sqrt{n}}{n^{1/6}} \le x\biggr)
&= F_1(x), \\
\lim_{n\to\infty} 
\Prob\biggl( \frac{\tilde{L}_n^\symmu-2\sqrt{n}}{2^{2/3}n^{1/6}}\le x\biggr)
&= F_1(x)^2.
\end{align}
We also have convergence of all the moments.
\end{thm}

This result should be compared with the results on random permutation 
and random signed permutation where the limiting distribution 
was $F_2(x)$ and $F_2(x)^2$ under the same scaling of the above 
(see \cite{BDJ}, \cite{TW3}, \cite{Bo} and Theorem \ref{thm1} above).

\begin{rem}
  As in the permutation case, the length of the longest \emph{increasing} 
subsequence and the length of the longest \emph{decreasing} subsequence 
of random involutions have the same statistics. 
This can be seen by noting that there is a bijection 
(called Robinson-Schensted correspondence, see e.g. \cite{Kn}) 
between 
the set $\tilde{S}_n$ of involutions of $n$ letters 
and the set of standard Young tableaux 
of size $n$, and the rows and the columns of standard Young tableaux 
have the same statistics under the push forward of the uniform probability 
distribution on $\tilde{S}_n$ under this bijection.
\end{rem}


\subsection{$\beta$-Plancherel measure on the set of Young diagrams}

Let $Y_n$ be the set of Young diagrams, or equivalently partitions, 
of size $n$.
Given a partition $\lambda=(\lambda_1,\lambda_2,\cdots) \vdash n$, 
let $d_\lambda$ denote the number of standard Young tableaux of shape
$\lambda$. 
We introduce the $\beta$-Plancherel measure $M^\beta_n$ on $Y_n$ 
defined by 
\begin{equation}
   M^\beta_n(\lambda)=\frac{d_\lambda^\beta}
{\sum_{\mu\vdash n}d_\mu^\beta}, \qquad \lambda\in Y_n.
\end{equation}
When $\beta=2$, this is the Plancherel measure which arises in the
representation theory.
We are interest in the typical shape and the fluctuation of $\lambda$ 
where $\lambda$ is taken randomly from the probability space $Y_n$  
with $M^\beta_n$.

A motivation introducing the above measure is the result of 
Regev \cite{Re}. 
In \cite{Re}, it is proved that for fixed $\beta >0$ and fixed $l$,
as $n\to\infty$,
\begin{equation}\label{int-as-1}
   \sum_{\substack{\lambda\vdash n \\ \lambda_1\le l}}
\bigl(d_\lambda\bigr)^\beta
\sim \biggl[ \frac{l^{l^2/2}l^n}{(\sqrt{2\pi})^{(l-1)/2}n^{(l-1)(l+2)/4}}
\biggr]^\beta \frac{n^{(l-1)/2}}{l!}
\int_{\R^l} e^{-\frac12\beta l\sum_{j}x_j^2}
\prod_{j<k} |x_j-x_k|^\beta d^l x.
\end{equation}
The multiple integral on the right hand side is 
the Selberg integral which can be computed exactly for each $\beta$.
When $\beta=1,2,4$, this integral is 
the normalization constant of the probability 
density of eigenvalues in 
GOE, GUE and GSE, respectively (see e.g. \cite{Mehta}).
So the basic question is if the $\beta$ in the definition 
of the $\beta$-Plancherel measure corresponds to 
the $\beta$ in the random matrix theory.
%Using the results on random involutions and random permutations, 
%we find that this is the case when $\beta=1,2$.

The well-known Robinson-Schensted correspondence \cite{Sc} 
establishes a bijection between 
$S_n$ and the pairs of standard Young tableaux with the same
shape of size $n$, $\RS : \pi \mapsto (P(\pi),Q(\pi))$.
Especially, we obtain $\sum_{\mu\vdash n} d_\mu^2= |S_n|=n!$.
Moreover, under $\RS$, the length of the longest increasing subsequence 
of $\pi\in S_n$ is equal to the number of boxes in the first 
row of $P(\pi)$ (or equally of $Q(\pi)$).
Therefore under $\RS$, the Plancherel measure $M^2_n$ is 
simply the push forward of 
the uniform probability measure on $S_n$ to $Y_n$, 
and the number of boxes in the first row of a random Young diagram 
and the length of the longest increasing subsequence of a random 
permutation have the same statistics : GUE ($\beta=2$) 
fluctuation in the limit.

If $\RS(\pi)=(P,Q)$, then $\RS(\pi^{-1})=(Q,P)$ (see e.g. \cite{Kn}).
Therefore the set of involutions $\tilde{S}_n$ is bijective to 
the set of (single) standard Young tableaux, and 
the number of boxes in the first row of a random Young diagram 
taken under the probability $M^1_n$ has the same statistics 
with the length of the longest increasing subsequence of a random
involution : GOE ($\beta=1$) fluctuation in the limit.

In fact, it is shown in \cite{Ok, BOO, kurtj:recent} that, 
for the case of $\beta=2$, in the large $n$ limit, 
the number of boxes in the $k^{\text{th}}$ of a random Young diagram 
corresponds to the $k^{\text{th}}$ largest eigenvalue 
of a random GUE matrix.
Also the typical shape of $\lambda$ is obtained in \cite{LS} and 
\cite{VK1} which is related to Wigner's semicircle law \cite{Kerov}.
On the other hand, the second row of a random Young diagram 
for $\beta=1$ is 
discussed in \cite{BR2} implying that it corresponds to 
the second eigenvalue of a random GOE matrix.
It would be interesting to obtain similar results for general row 
for $\beta=1$ and also for general $\beta$.

\subsection{Random turn vicious walker model}

Random permutations and random involutions arise also in certain random 
walk process in a one dimensional integer lattice.
We call a particle left-movable (resp. right-movable) 
if its left (resp. right) site is vacant.
Initially, $p$ particles are located at the points $1,2,3,\cdots,p$.
At each time step $t$ ($t=1,2,\cdots$), 
one particle 
among the left-movable particles 
is selected at random,  
and is moved to its left site
(so at $t=1$, the leftmost particle is moved.)
Suppose this process is repeated for $n$ time steps. 
It is found that \cite{forrester99}
there is a bijection between all the possible configurations 
and the set of Young tableaux with at most $p$ rows of size $n$ : 
simply, in the $k^{\text{th}}$ row, write the times 
when the particle is originally at the point $k$ (see Figure \ref{fig-walk}).
\begin{figure}[ht]
 \centerline{\epsfig{file=walk.eps, width=13cm}}
 \caption{random turn vicious walkers and Young tableaux}
\label{fig-walk}
\end{figure}
In this correspondence the number of moves made by the $k^{\text{th}}$ 
particle (counted from the left) is the number of boxes 
in the $k^{\text{th}}$ row of the corresponding Young tableau.
If we take the limit $p\to\infty$, we remove the constraint 
on the number of rows. 
In other words, we start with countably many particles located at 
$1,2,3,\cdots$ and at each 
time step, we move a randomly selected left-movable particle 
to its left site.
Then the statistics of the number of moves made by the leftmost particle 
in $n$ time steps 
is identical to that of the number of boxes 
in the first row of a random Young 
diagram under the probability $M^1_n$ (or the length of the longest 
increasing subsequence of a random involution).
Hence in the large $n$ limit, we obtain GOE fluctuation.

Now suppose after $n$ steps of left moves, 
by taking $n$ steps of right moves, 
we want the particles brought back to their original positions.
The first $n$ steps gives a Young tableaux and the next $n$ step gives another 
Young tableaux.
So we are in the situation of pairs of Young tableaux, i.e. $\beta=2$ case.
Especially if the number of particles are infinite, 
the statistics of the moves made by the leftmost 
particle is identical to the length of the longest increasing 
subsequence of a random permutation : GUE fluctuation
\cite{forrester99}.






\subsection{Symmetrized versions of Johansson's model}\label{subsec-Jo}

In \cite{kurtj:shape}, Johansson introduced a certain model 
which has several probabilistic interpretations, 
as a randomly growing Young diagram, a totally asymmetric one dimensional 
exclusion process, a certain zero-temperature directed polymer in a random 
environment or as a kind of first-passage site percolation model.
Here we consider the symmetrized versions.
The limiting distributions are parallel to the results of Section 
\ref{sec-result} depending on the symmetry type.

Each model is defined on $\M$, a subset of $\Z^2$, and at each site 
$(i,j)\in\M$, we define a random variable $w(i,j)$ of geometric
distribution. 
We are interested in ``the length of the longest up/right path''.
Let $g(q)$ denote the geometric distribution with parameter $q$.
Let $0<q<1$, and let $\alpha,\beta\ge 0$ such that 
$\alpha \sqrt{q}, \beta \sqrt{q}<1$.
Otherwise stated, the random variables $w(i,j)$ are independent each other.
We denote by $(j,k)\nearrow (j',k')$ the set of up/right paths $\pi$ from 
$(j,k)$ to $(j',k')$.

\begin{enumerate}
\item Let $\M=\Z_+^2$. 
Let 
\begin{equation}
   w(i,j)\sim g(q).
\end{equation}
Define 
\begin{equation}
  G^\symmU(N)=\max \{\sum_{(i,j)\in\pi} w(i,j) : \pi\in (1,1)\nearrow(N,N) \}.
\end{equation}
\item Let $\M=\Z^2$. 
Let $w(i,j)=0$ if $i=0$ or $j=0$, and for $i,j\neq 0$, let 
\begin{equation}
   w(i,j)=w(-i,-j)\sim g(q).
\end{equation}
Define 
\begin{equation}
   G^\symmUU(N)=\max \{ \sum_{(i,j)\in\pi} w(i,j) : 
\pi\in (-N,-N)\nearrow(N,N) \}.
\end{equation}
\item Let $\M=\Z_+^2$. 
Let 
\begin{eqnarray}
   w(i,j)=w(j,i) &\sim& g(q),\quad   i\neq j, \\
   w(i,i)&\sim& g(\alpha\sqrt{q}).
\end{eqnarray}
Define 
\begin{equation}
  G^\symmO(N)=\max \{ \sum_{(i,j)\in\pi} w(i,j) : \pi\in (1,1)\nearrow(N,N) \}.
\end{equation}
\item Let $\M=\Z_+\times \Z_-$.
Let 
\begin{eqnarray}
  w(i,-j)=w(j,-i) &\sim& g(q), \quad i\neq -j, \\
  w(i,-i) &\sim& g(\beta\sqrt{q}).
\end{eqnarray}
Define 
\begin{equation}
   G^\symmS(N)=\max \{ \sum_{(i,j)\in\pi} w(i,j) : 
\pi\in (1,-N)\nearrow (N,-1) \}.
\end{equation}
\item Let $\M=\Z^2$.
Let $w(i,j)=0$ if $i=0$ or $j=0$. Otherwise 
\begin{eqnarray}
   w(i,j)=w(-i,-j) &\sim& g(q),\quad |i|\neq|j|, \\
   w(i,i)=w(-i,-i) &\sim& g(\alpha\sqrt{q}), \\
   w(i,-i)=w(-i,i) &\sim& g(\beta\sqrt{q}).
\end{eqnarray}
Define 
\begin{equation}
   G^\symmu(N)=\max \{ \sum_{(i,j)\in\pi} w(i,j) : 
\pi\in(-N,-N)\nearrow (N,N) \}.
\end{equation}
\end{enumerate}

We are interested in the limiting distributions of $G^\symmg(N)$.
\begin{thm}\label{thm-4.2}
Set 
\begin{equation}
  \eta(q)=\frac{2\sqrt{q}}{1-\sqrt{q}}, 
\qquad \rho(q)=\frac{q^{1/6}(1+\sqrt{q})^{1/3}}{1-\sqrt{q}}.
\end{equation}
We have 
\begin{enumerate}
\item 
\begin{equation}\label{e-appl-2}
  \lim_{N\to\infty}
\Prob\biggl(\frac{G^\symmU(N;q)-\eta(q)N}
{\rho(q)N^{1/3}}\le x\biggr)
= F_2(x).
\end{equation}
\item 
\begin{equation}
  \lim_{N\to\infty}
\Prob\biggl(\frac{G^\symmUU(N;q)-\eta(q)(2N)}
{2^{2/3}\rho(q)(2N)^{1/3}}\le x\biggr)
= F_2(x)^2.
\end{equation}
\item 
\begin{eqnarray}
  \lim_{N\to\infty}
\Prob\biggl(\frac{G^\symmO(N;q,\alpha)-\eta(q)N}
{\rho(q)N^{1/3}}\le x\biggr)
&=& F_4(x), \qquad 0\le\alpha<1, \\
\label{e-appl-3}
  \lim_{N\to\infty}
\Prob\biggl(\frac{G^\symmO(N;q,\alpha)-\eta(q)N}
{\rho(q)N^{1/3}}\le x\biggr)
&=& F^\symmO(x;w), \quad \alpha=1-\frac{2w}{\rho(q)N^{1/3}}, \\
  \lim_{N\to\infty}
\Prob\biggl(\frac{G^\symmO(N;q,\alpha)-\eta(q)N}
{\rho(q)N^{1/3}}\le x\biggr)
&=& 0, \qquad\qquad \alpha>1.
\end{eqnarray}
\item
\begin{equation}
  \lim_{N\to\infty}
\Prob\biggl(\frac{G^\symmS(N;q,\beta)-\eta(q)N}
{\rho(q)N^{1/3}}\le x\biggr)
= F_1(x), \qquad 0\le\beta.
\end{equation}
\item 
\begin{eqnarray}
  \lim_{N\to\infty}
\Prob\biggl(\frac{G^\symmu(N;q,\alpha,\beta)-\eta(q)(2N)}
{2^{2/3}\rho(q)(2N)^{1/3}}\le x\biggr)
&=& F_2(x), \qquad 0\le\alpha<1, \beta\ge 0, \\
  \lim_{N\to\infty}
\Prob\biggl(\frac{G^\symmu(N;q,\alpha,\beta)-\eta(q)(2N)}
{2^{2/3}\rho(q)(2N)^{1/3}}\le x\biggr)
&=& F^\symmu(x;w), \quad \alpha=1-\frac{2w}{\rho(q)(2N)^{1/3}}, \beta\ge 0 \\
  \lim_{N\to\infty}
\Prob\biggl(\frac{G^\symmu(N;q,\alpha,\beta)-\eta(q)(2N)}
{2^{2/3}\rho(q)(2N)^{1/3}}\le x\biggr)
&=& 0, \qquad\qquad \alpha>1, \beta\ge 0.
\end{eqnarray}
\end{enumerate}
\end{thm}

\begin{rem}
  The results \eqref{e-appl-2} and \eqref{e-appl-3} with $w=0$ 
are obtained in \cite{kurtj:recent}. 
For $\tsymmU$, 
the longest up/right path in $(1,1)\nearrow (M,N)$, $M\neq N$, 
is also considered.
\end{rem}

The proof of the above theorem is analogous to that of 
the result in Section \ref{sec-result}.
There again is a Toeplitz/Hankel determinantal formula for each 
case. 
In fact, the only change with the symmetrized permutation case is 
that we have $(1+q+2\sqrt{q}\cos\theta)^N d\theta/(2\pi)$ instead of 
$e^{2t\cos\theta} d\theta/(2\pi)$ for the weight. 
In analyzing the asymptotics of orthogonal polynomials, we need
different scaling, but once we scale, the analysis is parallel.
One important common property of the above two weights is 
that both of the supports of their equilibrium measures change 
from the full circle to a part of the circle with one gap 
depending on $2t/k<1$ and $2t/k>1$, and $N/k<\eta(q)^{-1}$ 
and $N/k>\eta(q)^{-1}$. 
And also in the one gap case, 
the equilibrium measures decay like a square root at the ends 
of the supports.

